\section*{Abstract}
% wordcount 147
Sensory perception originates from the responses of sensory neurons, which react to a collection of sensory signals linked to various physical attributes of a singular perceptual object. 
Unraveling how the brain extracts perceptual information from these neuronal responses is a pivotal challenge in both computational neuroscience and machine learning.
Here we introduce a statistical mechanical theory, where perceptual information is first encoded in the correlated variability of sensory neurons and then reformatted into the firing rates of downstream neurons.
Applying this theory, we illustrate the encoding of motion direction using neural covariance and demonstrate high-fidelity direction recovery by spiking neural networks.
Networks trained under this theory also show enhanced performance in classifying natural images, achieving higher accuracy and faster inference speed. 
Our results challenge the traditional view of neural covariance as a secondary factor in neural coding, highlighting its potential influence on brain function.